%=============================================================================!
% 12.0  CHAPTER OPENING                                                       !
%=============================================================================!
A trajectory records where the atoms move, but the quantities we seek are
usually averages: the temperature of a liquid, its mean density or the
rate at which an atom hops between sites. To interpret these averages we
need statistical mechanics, which describes the probabilities of the
states available to a system. We begin with probability and ask when an
average along one trajectory can represent an average over those states.
Examples for which it cannot make the limitation explicit.

Counting the states of an isolated system then gives definitions of
entropy and temperature. Allowing the system to exchange energy with its
surroundings leads to the canonical distribution, from which we obtain
the velocity distribution and the sharing of kinetic energy among the
degrees of freedom counted in Chapters~5 and~11. Liquid-argon runs let us
compare these predictions with measured temperatures and fluctuations.
Chapter~13 uses the resulting distributions to test thermostats: holding
the mean temperature is only one part of sampling the required ensemble.
